\begin{afterword}
\summaryhead

The post begins with semantic priming. People who have just seen the word ``water'' recognize ``drink'' as a word faster. Priming acts at the level of recognizing letters, spreads to many associated words, and is, the author says, subconscious and unstoppable; the Stroop test is offered as a demonstration.

Anchors work partly the same way. People asked whether the mean temperature in Germany is above or below 5°C were then faster to recognize words like ``cold'' and ``snow,'' and people asked about 20°C were faster with ``hot'' and ``sun.'' The wider finding, that uninformative, false or irrelevant information affects estimates, is called contamination. Adjustment from an anchor still explains some cases, in particular anchors people produce themselves, but the author says modern research attributes most anchoring to contamination. Purchase-limit signs in shops, the post adds, have been shown to work.

Finally, contamination is called a form of confirmation bias built into the brain. An idea primes thoughts compatible with it and so, the post says, ensures its own continued existence, and one quick flash can decide the conclusion.

\responsehead

Most of the post is sound reporting. The priming it describes, faster recognition of related words, is well established. The temperature experiment is reported correctly. The post also brings its account of anchoring up to date: a month earlier ``Anchoring and Adjustment'' called adjustment the current theory, and this post says contamination explains most anchoring, in line with the consensus its source describes. The claim that priming is ``unstoppable'' is stated as fact, although researchers had argued before the post that semantic processing depends on attention.

The problem is the last two paragraphs. They move from how fast people recognize words to how people hold conclusions. The link to confirmation bias has a basis: in the account of anchoring the post relies on, people test whether the anchor is right, and tests of that kind tend to seek confirmation. But the post goes further. It says a primed idea ``thereby ensures its continued existence,'' that confirmation bias is ``built directly into our hardware,'' and that after ``one quick flash'' the ``bottom line is already decided.'' Every study in the post measures recognition times or estimates given in the same session. None measures whether a primed idea lasts, or whether it decides anything later.

Later research bears on the difference between the two parts. Word-recognition priming, the post's first topic, has been replicated. After the post, several findings that brief primes change later judgments and behavior failed to replicate. The post's conclusion is a claim of that second kind, though not one of the claims that failed, and the post could not have known of the failures.

In short: an accurate account of word priming and of one anchoring experiment, extended in the last two paragraphs to a claim about how beliefs last that none of its evidence tests.
\end{afterword}
